\chapter{Code Style Guide}

This appendix defines the coding standards and conventions used throughout the book. All code examples follow these guidelines to ensure consistency and readability.

\section{Python Style Standards}

\subsection{General Principles}

All Python code in this book adheres to:

\begin{itemize}
    \item \textbf{PEP 8} style guide for Python code
    \item \textbf{PEP 257} for docstring conventions
    \item \textbf{Type hints} (PEP 484) for function signatures
    \item \textbf{Black} code formatter (line length: 88 characters)
    \item \textbf{isort} for import sorting
\end{itemize}

\subsection{Naming Conventions}

\begin{lstlisting}[style=python, caption=Naming Convention Examples]
# Classes: PascalCase
class FeatureStore:
    pass

class ModelTrainer:
    pass

# Functions and variables: snake_case
def train_model(data: pd.DataFrame) -> Model:
    model_name = "fraud_detector"
    training_data = prepare_data(data)
    return model

# Constants: UPPER_SNAKE_CASE
MAX_RETRIES = 3
DEFAULT_TIMEOUT = 30
API_BASE_URL = "https://api.example.com"

# Private methods/attributes: leading underscore
class DataProcessor:
    def __init__(self):
        self._cache = {}

    def _validate_data(self, data):
        pass

# Type aliases: PascalCase
FeatureDict = Dict[str, Union[int, float, str]]
ModelConfig = Dict[str, Any]
\end{lstlisting}

\subsection{Import Organization}

Imports are organized in three groups, separated by blank lines:

\begin{lstlisting}[style=python, caption=Import Organization]
# Standard library imports
import os
import sys
from datetime import datetime
from pathlib import Path
from typing import Dict, List, Optional, Any, Tuple

# Third-party imports
import numpy as np
import pandas as pd
from sklearn.ensemble import RandomForestClassifier
import mlflow
import mlflow.sklearn

# Local application imports
from src.data.loaders import DataLoader
from src.features.transformers import FeatureTransformer
from src.models.trainer import ModelTrainer
from src.utils.logger import setup_logger
\end{lstlisting}

\subsection{Type Hints}

All function signatures include type hints:

\begin{lstlisting}[style=python, caption=Type Hint Examples]
from typing import Dict, List, Optional, Union, Tuple, Any
import pandas as pd
import numpy as np

# Simple types
def calculate_score(value: float) -> float:
    return value * 2.0

# Complex types
def process_features(
    data: pd.DataFrame,
    feature_names: List[str],
    config: Dict[str, Any]
) -> Tuple[np.ndarray, List[str]]:
    """Process features with given configuration."""
    processed_data = data[feature_names].values
    return processed_data, feature_names

# Optional types
def load_model(path: str, version: Optional[str] = None) -> Model:
    """Load model from path, optionally specific version."""
    if version is None:
        version = "latest"
    return Model.load(path, version)

# Union types
def parse_value(value: Union[str, int, float]) -> float:
    """Parse value to float."""
    return float(value)

# Generic types
from typing import TypeVar, Generic

T = TypeVar('T')

class Cache(Generic[T]):
    def __init__(self) -> None:
        self._data: Dict[str, T] = {}

    def get(self, key: str) -> Optional[T]:
        return self._data.get(key)

    def set(self, key: str, value: T) -> None:
        self._data[key] = value
\end{lstlisting}

\subsection{Docstrings}

Use Google-style docstrings for all public functions and classes:

\begin{lstlisting}[style=python, caption=Docstring Examples]
def train_model(
    X_train: np.ndarray,
    y_train: np.ndarray,
    model_config: Dict[str, Any],
    validation_split: float = 0.2
) -> Tuple[Model, Dict[str, float]]:
    """Train machine learning model with cross-validation.

    This function trains a model using the provided training data
    and configuration. It performs internal validation and returns
    both the trained model and validation metrics.

    Args:
        X_train: Training features array of shape (n_samples, n_features)
        y_train: Training labels array of shape (n_samples,)
        model_config: Configuration dictionary containing model parameters
        validation_split: Fraction of data to use for validation (default: 0.2)

    Returns:
        Tuple containing:
            - Trained model instance
            - Dictionary of validation metrics (accuracy, precision, recall, f1)

    Raises:
        ValueError: If validation_split is not in range [0, 1]
        RuntimeError: If model training fails

    Example:
        >>> model_config = {"n_estimators": 100, "max_depth": 10}
        >>> model, metrics = train_model(X_train, y_train, model_config)
        >>> print(f"Accuracy: {metrics['accuracy']:.3f}")
        Accuracy: 0.923

    Note:
        The model is automatically logged to MLflow if tracking is enabled.
        See Chapter 8 for details on experiment tracking.

    References:
        See Section 7.3 for model architecture details.
    """
    if not (0 <= validation_split <= 1):
        raise ValueError("validation_split must be in range [0, 1]")

    # Implementation...
    pass

class FeatureStore:
    """Centralized feature storage and retrieval system.

    The FeatureStore provides a unified interface for storing and
    retrieving features for machine learning models. It supports
    both batch and real-time feature access.

    Attributes:
        connection: Database connection instance
        cache: In-memory feature cache
        feature_specs: Dictionary of feature specifications

    Example:
        >>> store = FeatureStore(connection_string="postgresql://...")
        >>> features = store.get_features(user_id="123", feature_names=["age", "income"])
        >>> print(features)
        {"age": 35, "income": 75000}

    See Also:
        Chapter 6: Feature Engineering and Feature Stores
    """

    def __init__(self, connection_string: str) -> None:
        """Initialize feature store connection.

        Args:
            connection_string: Database connection string
        """
        self.connection = self._create_connection(connection_string)
        self.cache = {}
        self.feature_specs = {}
\end{lstlisting}

\subsection{Error Handling}

Explicit error handling with informative messages:

\begin{lstlisting}[style=python, caption=Error Handling Patterns]
import logging
from typing import Optional

logger = logging.getLogger(__name__)

# Specific exceptions
class ModelNotFoundError(Exception):
    """Raised when model cannot be found in registry."""
    pass

class DataValidationError(Exception):
    """Raised when data validation fails."""
    pass

# Error handling with logging
def load_model(model_id: str) -> Model:
    """Load model with proper error handling."""
    try:
        logger.info(f"Loading model: {model_id}")
        model = Model.load(model_id)
        logger.info(f"Successfully loaded model: {model_id}")
        return model

    except FileNotFoundError as e:
        logger.error(f"Model file not found: {model_id}")
        raise ModelNotFoundError(
            f"Model {model_id} not found in registry"
        ) from e

    except Exception as e:
        logger.exception(f"Unexpected error loading model: {model_id}")
        raise

# Validation with early returns
def validate_config(config: Dict[str, Any]) -> None:
    """Validate configuration dictionary."""
    if not config:
        raise ValueError("Configuration cannot be empty")

    if "model_type" not in config:
        raise ValueError("Configuration must specify 'model_type'")

    if config["model_type"] not in ["classifier", "regressor"]:
        raise ValueError(
            f"Invalid model_type: {config['model_type']}. "
            f"Must be 'classifier' or 'regressor'"
        )

    logger.info("Configuration validated successfully")
\end{lstlisting}

\section{Configuration Management}

\subsection{YAML Configuration Files}

All configuration uses YAML format:

\begin{lstlisting}[style=yaml, caption=Configuration File Structure]
# config/training_config.yaml
model:
  name: "fraud_detector"
  type: "xgboost"
  hyperparameters:
    n_estimators: 100
    max_depth: 10
    learning_rate: 0.1
    subsample: 0.8

data:
  training_path: "s3://bucket/data/train.parquet"
  validation_path: "s3://bucket/data/val.parquet"
  test_path: "s3://bucket/data/test.parquet"
  batch_size: 1000

training:
  num_epochs: 10
  early_stopping_patience: 3
  validation_split: 0.2
  random_seed: 42

logging:
  level: "INFO"
  log_dir: "logs/"
  mlflow_tracking_uri: "http://mlflow:5000"

deployment:
  serving_endpoint: "http://api:8000/predict"
  max_batch_size: 100
  timeout_ms: 100
\end{lstlisting}

\subsection{Configuration Loading}

\begin{lstlisting}[style=python, caption=Configuration Loading Pattern]
from pathlib import Path
from typing import Dict, Any
import yaml
from dataclasses import dataclass

@dataclass
class ModelConfig:
    """Model configuration data class."""
    name: str
    type: str
    hyperparameters: Dict[str, Any]

    @classmethod
    def from_dict(cls, config_dict: Dict[str, Any]) -> 'ModelConfig':
        """Create ModelConfig from dictionary."""
        return cls(**config_dict)

def load_config(config_path: str) -> Dict[str, Any]:
    """Load configuration from YAML file.

    Args:
        config_path: Path to YAML configuration file

    Returns:
        Dictionary containing configuration

    Raises:
        FileNotFoundError: If config file doesn't exist
        ValueError: If YAML is invalid
    """
    path = Path(config_path)

    if not path.exists():
        raise FileNotFoundError(f"Config file not found: {config_path}")

    try:
        with open(path, 'r') as f:
            config = yaml.safe_load(f)
        return config
    except yaml.YAMLError as e:
        raise ValueError(f"Invalid YAML in config file: {e}")

# Usage
config = load_config("config/training_config.yaml")
model_config = ModelConfig.from_dict(config['model'])
\end{lstlisting}

\section{Logging Standards}

\subsection{Logger Setup}

\begin{lstlisting}[style=python, caption=Logging Configuration]
import logging
import sys
from pathlib import Path

def setup_logger(
    name: str,
    level: str = "INFO",
    log_file: Optional[str] = None
) -> logging.Logger:
    """Configure logger with consistent format.

    Args:
        name: Logger name (usually __name__)
        level: Logging level (DEBUG, INFO, WARNING, ERROR, CRITICAL)
        log_file: Optional file path for logging

    Returns:
        Configured logger instance
    """
    logger = logging.getLogger(name)
    logger.setLevel(getattr(logging, level.upper()))

    # Console handler
    console_handler = logging.StreamHandler(sys.stdout)
    console_handler.setLevel(logging.INFO)

    # Formatter
    formatter = logging.Formatter(
        fmt='%(asctime)s - %(name)s - %(levelname)s - %(message)s',
        datefmt='%Y-%m-%d %H:%M:%S'
    )
    console_handler.setFormatter(formatter)
    logger.addHandler(console_handler)

    # File handler (optional)
    if log_file:
        file_handler = logging.FileHandler(log_file)
        file_handler.setLevel(logging.DEBUG)
        file_handler.setFormatter(formatter)
        logger.addHandler(file_handler)

    return logger

# Usage in modules
logger = setup_logger(__name__)

logger.debug("Detailed debugging information")
logger.info("General information messages")
logger.warning("Warning messages")
logger.error("Error messages")
logger.critical("Critical error messages")
\end{lstlisting}

\section{Testing Patterns}

\subsection{Unit Test Structure}

\begin{lstlisting}[style=python, caption=Unit Test Example]
import pytest
import numpy as np
import pandas as pd
from src.models.trainer import ModelTrainer
from src.data.loaders import DataLoader

class TestModelTrainer:
    """Test suite for ModelTrainer class."""

    @pytest.fixture
    def sample_data(self):
        """Create sample training data."""
        np.random.seed(42)
        X = np.random.randn(100, 10)
        y = np.random.randint(0, 2, 100)
        return X, y

    @pytest.fixture
    def trainer(self):
        """Create ModelTrainer instance."""
        config = {
            "model": {
                "type": "random_forest",
                "hyperparameters": {
                    "n_estimators": 10,
                    "max_depth": 5
                }
            }
        }
        return ModelTrainer(config)

    def test_train_model_shape(self, trainer, sample_data):
        """Test that trained model produces correct output shape."""
        X, y = sample_data
        model = trainer.train(X, y)

        predictions = model.predict(X)
        assert predictions.shape == y.shape

    def test_train_model_accuracy(self, trainer, sample_data):
        """Test that model achieves minimum accuracy."""
        X, y = sample_data
        model = trainer.train(X, y)

        accuracy = model.score(X, y)
        assert accuracy > 0.5  # Better than random

    def test_invalid_config_raises_error(self):
        """Test that invalid configuration raises ValueError."""
        invalid_config = {"model": {"type": "invalid_model"}}

        with pytest.raises(ValueError):
            trainer = ModelTrainer(invalid_config)

    @pytest.mark.parametrize("n_estimators,expected_min_accuracy", [
        (10, 0.6),
        (50, 0.7),
        (100, 0.75),
    ])
    def test_accuracy_improves_with_trees(
        self,
        sample_data,
        n_estimators,
        expected_min_accuracy
    ):
        """Test that accuracy improves with more trees."""
        X, y = sample_data
        config = {
            "model": {
                "type": "random_forest",
                "hyperparameters": {"n_estimators": n_estimators}
            }
        }
        trainer = ModelTrainer(config)
        model = trainer.train(X, y)

        accuracy = model.score(X, y)
        assert accuracy >= expected_min_accuracy
\end{lstlisting}

\section{Code Comments}

\subsection{When to Comment}

\begin{enumerate}
    \item \textbf{Complex algorithms:} Explain non-obvious logic
    \item \textbf{Business rules:} Document domain-specific requirements
    \item \textbf{Workarounds:} Explain why hacky code is necessary
    \item \textbf{TODOs:} Mark incomplete implementations
    \item \textbf{Performance considerations:} Explain optimization choices
\end{enumerate}

\subsection{Comment Examples}

\begin{lstlisting}[style=python, caption=Good Comment Examples]
# Good: Explains WHY, not WHAT
def calculate_fraud_score(transaction: Dict) -> float:
    """Calculate fraud probability score."""
    # Use log transformation to handle skewed amount distribution
    # Reference: Section 5.2.3 for feature engineering rationale
    log_amount = np.log1p(transaction['amount'])

    # Apply higher weight to international transactions
    # based on historical fraud patterns (85% of fraud is international)
    int_multiplier = 2.5 if transaction['is_international'] else 1.0

    return log_amount * int_multiplier

# Bad: States the obvious
x = x + 1  # Increment x by 1

# Good: Documents business rule
MAX_TRANSACTION_AMOUNT = 10000  # Regulatory limit from PCI-DSS v3.2.1

# Good: Explains workaround
def process_batch(data):
    # HACK: pandas groupby has memory leak in version 1.2.3
    # Explicitly delete intermediate dataframes to avoid OOM
    # TODO: Remove after upgrading to pandas 1.3.0
    grouped = data.groupby('user_id')
    result = grouped.apply(aggregate_features)
    del grouped
    return result
\end{lstlisting}

\section{Code Organization Patterns}

\subsection{Module Structure}

Each module follows this structure:

\begin{lstlisting}[style=python, caption=Standard Module Structure]
"""Module docstring describing purpose.

This module provides functionality for...

Example:
    >>> from module import ClassName
    >>> instance = ClassName()
    >>> result = instance.method()
"""

# Imports (standard, third-party, local)
import os
from typing import Dict, List

import numpy as np

from .utils import helper_function

# Module-level constants
DEFAULT_CONFIG = {"param": "value"}

# Module-level variables (if necessary)
_cache: Dict[str, any] = {}

# Helper functions (private)
def _internal_helper(data: np.ndarray) -> np.ndarray:
    """Internal helper function."""
    pass

# Public functions
def public_function(param: str) -> str:
    """Public API function."""
    pass

# Classes
class PublicClass:
    """Main class implementation."""
    pass

# Main execution block
if __name__ == "__main__":
    # CLI or testing code
    pass
\end{lstlisting}

\section{Performance Patterns}

\subsection{Vectorization}

Prefer vectorized operations over loops:

\begin{lstlisting}[style=python, caption=Vectorization Examples]
import numpy as np
import pandas as pd

# Bad: Slow loop
result = []
for x in data:
    result.append(x ** 2 + 2 * x + 1)

# Good: Vectorized
result = data ** 2 + 2 * data + 1

# Bad: DataFrame iteration
for idx, row in df.iterrows():
    df.at[idx, 'result'] = row['a'] + row['b']

# Good: Vectorized DataFrame operations
df['result'] = df['a'] + df['b']

# Good: Using .apply() when vectorization isn't possible
df['result'] = df.apply(lambda row: complex_function(row), axis=1)
\end{lstlisting}

\section{Cross-Reference Style}

Use consistent cross-references throughout the book:

\begin{lstlisting}[language=TeX, caption=Cross-Reference Examples]
% Chapter references
See Chapter~\ref{ch:data_pipelines} for details on data processing.

% Section references
As discussed in Section~\ref{sec:feature_engineering}, features must be...

% Figure references
Figure~\ref{fig:architecture} shows the system architecture.

% Table references
Table~\ref{tab:metrics} compares model performance.

% Equation references
Using Equation~\eqref{eq:precision}, we calculate...

% Code listing references
Listing~\ref{lst:training} demonstrates the training loop.

% Multiple references
Chapters~\ref{ch:serving}--\ref{ch:monitoring} cover production operations.
\end{lstlisting}

\section{Summary}

These coding standards ensure consistency and quality across all book examples:

\begin{itemize}
    \item \textbf{Follow PEP 8} for all Python code
    \item \textbf{Use type hints} for function signatures
    \item \textbf{Write comprehensive docstrings} with examples
    \item \textbf{Handle errors explicitly} with informative messages
    \item \textbf{Use YAML for configuration} with validation
    \item \textbf{Set up proper logging} in all modules
    \item \textbf{Write tests} for critical functionality
    \item \textbf{Comment why, not what} for complex logic
    \item \textbf{Prefer vectorization} over loops
    \item \textbf{Use consistent cross-references} to other chapters
\end{itemize}

All code examples in this book are production-ready and can be adapted for real-world ML systems.
